\section{Approximate Rigidity}
\label{app:approx-rigidity}
The encoder characterization also gives a qualitative statement for approximate implementations when the admissible encoder class is compact. The common regularity and decoder bounds below are explicit additional assumptions.

Fix the desired function $f_{\phi,\psi}:I^K\to\R$ and impose the hypotheses of Theorem~\ref{thm:rigidity}. Let $\omega:[0,2]\to[0,\infty)$ satisfy $\omega(h)\to0$ as $h\downarrow0$. Let $\mathcal E_{P,\omega}$ consist of the encoders $E\in C(I;\R^2)$ satisfying
$$
 \|E\|_\infty\le\sqrt P,\quad
 |E(s)-E(t)|\le\omega(|s-t|)\quad\forall s,t\in I.
$$
Write $\mathcal A_\phi\subset C(I;\R^2)$ for the family $s\mapsto c+us+v\phi(s)$ with $c,u,v\in\R^2$ and linearly independent $u,v$, and use uniform distance to this family.

\begin{proposition}
\label{prop:approx-rigidity}
For fixed $P$, $\omega$, and $L>0$, there is a nondecreasing function $\eta:[0,\infty)\to[0,\infty)$ with $\lim_{\epsilon\downarrow0}\eta(\epsilon)=0$ and the following property. For every $\epsilon\ge0$, if $E\in\mathcal E_{P,\omega}$ and an $L$-Lipschitz decoder $D:\R^2\to\R$ satisfy
$$
 \sup_{\mathbf s\in I^K}
 |D(r_E(\mathbf s))-f_{\phi,\psi}(\mathbf s)|\le\epsilon,
$$
then $\operatorname{dist}_\infty(E,\mathcal A_\phi)\le\eta(\epsilon)$. The function $\eta$ may also depend on the fixed $K$, $\phi$, and $\psi$.
\end{proposition}
\begin{proof}
Approximate computation and decoder Lipschitz continuity imply, for every $\mathbf s,\mathbf t\in I^K$,
\begin{equation}
 |f(\mathbf s)-f(\mathbf t)|
 \le 2\epsilon+L|r_E(\mathbf s)-r_E(\mathbf t)|.
 \label{eq:approx-collision}
\end{equation}
The peak and common modulus bounds make $\mathcal E_{P,\omega}$ compact in the uniform norm by the Arzel\`a--Ascoli theorem. The class is closed because its defining inequalities survive uniform limits.

Define the continuous defect
$$
 d_L(E)=\frac12\sup_{\mathbf s,\mathbf t\in I^K}
 \bigl[|f(\mathbf s)-f(\mathbf t)|
       -L|r_E(\mathbf s)-r_E(\mathbf t)|\bigr]_+.
$$
Indeed, $|d_L(E)-d_L(\widetilde E)|\le KL\|E-\widetilde E\|_\infty$. Set $\eta(\epsilon)$ to the supremum of $\operatorname{dist}_\infty(E,\mathcal A_\phi)$ over $E\in\mathcal E_{P,\omega}$ with $d_L(E)\le\epsilon$, taking the supremum to be zero when the set is empty. The distance is continuous even though independence makes $\mathcal A_\phi$ nonclosed. Compactness makes this supremum finite, and \eqref{eq:approx-collision} proves the required bound.

If $\eta(\epsilon)$ failed to tend to zero, a sequence $\epsilon_n\downarrow0$ would supply encoders at a fixed positive distance from $\mathcal A_\phi$. A subsequence converges uniformly to $E\in\mathcal E_{P,\omega}$. Passing to the limit gives
$$
 |f(\mathbf s)-f(\mathbf t)|
 \le L|r_E(\mathbf s)-r_E(\mathbf t)|.
$$
Thus $f$ is identifiable through $E$, and Theorem~\ref{thm:rigidity} places $E$ in $\mathcal A_\phi$ with linearly independent coefficient vectors $u,v\in\R^2$. Uniform convergence then contradicts the positive distance. This proof requires no convergence of the decoders.
\end{proof}

This proposition supplies neither an explicit rate nor a quantitative parameter-identification guarantee. Its assumptions hold, for example, for encoders with a prescribed common Lipschitz constant and decoders with a fixed disturbance-amplification budget. Peak power alone does not supply the compactness used in the proof. We do not assert that these additional assumptions are necessary or that the proposition resolves quantitative approximate rigidity.
